\subsection{群 X(H)}

设 H 是紧李群。用 X(H) 表示从 H 到拓扑群 \(C^{*}\) 的连续同态所成的（交换）群。由第 III 章 §8 第 1 节定理 1，X(H) 的元素都是李群态射；对一切 \(a \in \mathrm{X}(\mathrm{H})\) ，a 的微分是一个 R-线性映射 \(\mathrm{L}(a) : \mathrm{L}(\mathrm{H}) \to \mathrm{L}(\mathbf{C}^{*})\) 。从现在起，我们把 \(C^{*}\) 的李代数与 C 这样等同起来，使得 \(C^{*}\) 的指数映射与映射 \(z \mapsto e^{z}\) （从 C 到 \(C^{*}\) ）相合。于是，每个元素 \(a \in \mathrm{X}(\mathrm{H})\) 对应一个元素 \(\mathrm{L}(a) \in \mathrm{Hom}_{\mathbf{R}}(\mathrm{L}(\mathrm{H}), \mathbf{C})\) ；我们用 \(\delta(a)\) 表示与之对应的 \(\mathrm{Hom}_{\mathbf{C}}(\mathrm{L}(\mathrm{H})_{(\mathbf{C})}, \mathbf{C})\) 中的元素（即它在 \(\mathrm{L}(\mathrm{H}) \subset \mathrm{L}(\mathrm{H})_{(\mathbf{C})}\) 上的限制等于 \(\mathrm{L}(a)\) ）。

对一切 \(x \in \mathrm{L}(\mathrm{H})\) 与一切 \(a \in \mathrm{X}(\mathrm{H})\) ，我们有

\[
a (\exp_ {\mathrm{H}} x) = e ^ {\delta (a) (x)},
\]

这由指数映射的函子性得出（第 III 章，§6，第 4 节，命题 10）。

我们常把群 \(\mathrm{X}(\mathrm{H})\) 写成加法；这时，我们把元素 \(a(g)\) （\(C^{*}\) 中的）记作 \(g^{a}\) 。采用这一记号，我们有公式

\[
g ^ {a + b} = g ^ {a} g ^ {b}, \quad g \in \mathrm{H}, \quad a, b \in \mathrm{X(H)},
\]

以及

\[
(\exp_ {\mathrm{H}} x) ^ {a} = e ^ {\delta (a) (x)}, x \in \mathrm{L} (\mathrm{H}), a \in \mathrm{X} (\mathrm{H}).
\]

由于 H 紧，X(H) 的元素取值于模为 1 的复数组成的子群 \(\mathbf{U} = \mathbf{U}(1, \mathbf{C})\) 中，故 X(H) 可与从 H 到 U 的连续（或解析）同态所成的群等同起来。由此得出，对一切 \(a \in \mathrm{L}(\mathrm{H})\) ，映射 \(\mathrm{L}(a)\) 取值于 C 的子空间 Ri 中，故 \(\delta(a)\) 把 \(\mathrm{L}(\mathrm{H})\) 映到 Ri 中。

若 H 交换，则 \(X(H)\) 就是 H 的（离散）对偶群（《谱理论》，第 II 章，§1，第 1 节）。若 H 交换且有限，则 \(X(H)\) 可与对偶有限群 \(D(H) = \text{Hom}_{\mathbf{Z}}(H, \mathbf{Q}/\mathbf{Z})\) 等同起来（其中如同《代数学》第 VII 章 §4 第 9 节例 1 那样，把 Q/Z 与 \(C^{*}\) 的一个子群借助同态 \(r \mapsto \exp(2\pi ir)\) 等同起来）。

对紧李群的任一态射 \(f: \mathrm{H} \to \mathrm{H}'\) ，我们用 \(\mathrm{X}(f)\) 表示同态 \(a \mapsto a \circ f\) （从 \(\mathrm{X}(\mathrm{H}')\) 到 \(\mathrm{X}(\mathrm{H})\) ） 。若 \(\mathrm{K}\) 是紧李群 H 的闭正规子群，则有 Z-模正合列 \(0 \rightarrow X(H/K) \rightarrow X(H) \rightarrow X(K)\) 。

命题 1. 对每个紧李群 H，Z-模 X(H) 是有限型的。当 H 连通时它是自由的。

先假定 H 连通；X(H) 的每个元素在 H 的换位子群 D(H) 上取值零，故我们有同态 \(\mathrm{X}(\mathrm{H}/\mathrm{D}(\mathrm{H}))\to\mathrm{X}(\mathrm{H})\) 。而 \(\mathrm{H}/\mathrm{D}(\mathrm{H})\) 连通且交换，故是环面，且 \(\mathrm{X}(\mathrm{H}/\mathrm{D}(\mathrm{H}))\) 是有限型自由 Z-模（《谱理论》，第 II 章，§2，第 1 节，命题 1 的推论 2）。在一般情形，由正合列

\[
0 \to \mathrm{X} (\mathrm{H} / \mathrm{H} _ {0}) \to \mathrm{X} (\mathrm{H}) \to \mathrm{X} (\mathrm{H} _ {0}),
\]

其中 \(X(H_{0})\) 有限型自由且 \(X(H/H_{0})\) 有限，得出 \(X(H)\) 是有限型的。

命题 2. 设 H 是交换紧李群，\((a_{i})_{i\in I}\) 是 X(H) 中元素的一个族；诸 \(a_{i}\) 生成 X(H) 当且仅当诸核 Ker \(a_{i}\) 之交约化为单位元。

由《谱理论》第 II 章 §1 第 7 节定理 4，\(a_i\) 的核的正交补是子群 \(A_i\) （\(X(H)\) 中由 \(a_i\) 生成的）；由同前定理 4 的推论 2，\(\bigcap \operatorname{Ker} a_i\) 的正交补是 \(X(H)\) 中由诸 \(A_i\) 生成的子群，命题得证。

